% Round 16. Separate continuation after the principal results were frozen.
% Root and global_height independent full reviews PASS at mathematical SHA
% e8bc3c712d9ea91d36d4c9b10efb7b38c9a37fb7545c83ad20e5f59ad04f2d5f.
% This activation changes only the status notice.
% Subsequent PDF layout only: GG16.5 is displayed on two lines.
% Reduced models only; no priority claim.
\section{A graph correction for scalar total boundary images}
\label{r16gg:section}

\paragraph{Provenance and exact scope.}
Historical Round 7 already studied scalar-degree correction matrices
and recorded higher-codimension numerical admissibility criteria in
\path{history_01_07/bb_round7/proofs/threefold_scalar_degree_kernel.tex}.
Its original review flags remain unchanged. The result here is a
direct sufficient condition on the \emph{total cohomological Gysin
image at each component}, followed by native actual activation.
It neither reasserts all historical higher-kernel conclusions nor
claims priority for the graph method. The general-degree support
argument below is proved explicitly; the already reviewed
Section~\ref{r16mb:section} states its degree-four case only.

\begin{lemma}[Single-divisor generators of the supported image in every degree]
\label{r16gg:support}
Let $D/k$ be a smooth projective $n$-fold over a finite field,
let $B=\bigcup_a E_a$ be a reduced strict normal crossings divisor
whose intersections are smooth of their expected codimensions,
and let $\ell\ne\operatorname{char}k$. For every $1\leq r\leq n$,
\[
 \operatorname{im}\left(H^{2r}_{B_{\bar k}}(D_{\bar k},\mathbf Q_\ell)
                   \longrightarrow H^{2r}(D_{\bar k},\mathbf Q_\ell)\right)
 =\sum_a\operatorname{im}\left(
 H^{2r-2}((E_a)_{\bar k},\mathbf Q_\ell)(-1)
                  \xrightarrow{j_{a*}}H^{2r}(D_{\bar k},\mathbf Q_\ell)
                         \right).
 \tag{GG16.1}\label{r16gg:support-image}
\]
The statement is also valid for an empty boundary, with zero image.
\end{lemma}

\begin{proof}
After a finite residue extension all geometric components can be
labeled separately; this changes neither geometric cohomology nor
the claimed images. Iterating the closed-support Mayer--Vietoris
triangle gives the bounded Frobenius-equivariant spectral sequence
\[
 E_1^{-p,q}=\bigoplus_{|I|=p+1}H^q_{(E_I)_{\bar k}}
                  (D_{\bar k},\mathbf Q_\ell)
       \Longrightarrow H^{q-p}_{B_{\bar k}}(D_{\bar k},\mathbf Q_\ell),
                         \qquad p\geq0.
 \tag{GG16.2}\label{r16gg:support-sequence}
\]
Here $E_I=\bigcap_{a\in I}E_a$; empty intersections are omitted.
This is the resolution for closed supports, not ordinary cohomology
of the boundary union. Smooth-pair purity, with $t=p+1$, identifies
its terms as
\[
 H^q_{(E_I)_{\bar k}}(D_{\bar k},\mathbf Q_\ell)
                =H^{q-2t}((E_I)_{\bar k},\mathbf Q_\ell)(-t).
 \tag{GG16.3}\label{r16gg:purity}
\]
Every stratum is smooth and projective. Weil purity therefore
makes this term pure of weight $q$. The differential $d_1$
preserves $q$. Every $E_2^{-p,q}$ is consequently pure of
weight $q$, and each $d_s$ for $s\geq2$ changes the second
index to $q-s+1$. Disjoint Weil weights force these later
differentials to vanish, so $E_2=E_\infty$.

In total degree $2r$ one has $q=2r+p$, and the corresponding
stratum cohomology group is
\[
       H^{2r-p-2}((E_I)_{\bar k},\mathbf Q_\ell)(-p-1),
                    \qquad\text{of weight }2r+p.
 \tag{GG16.4}\label{r16gg:indices}
\]
Only
\[
       0\leq p\leq\min(n-1,\,2r-2,\,2n-2r)
\]
can occur: the stratum has dimension $n-p-1$, and its
cohomological degree lies between zero and twice that dimension.
Thus the unique possible weight $2r$ comes from column zero,
and every other graded piece has strictly larger weight.

Column zero is a subcomplex of the finite support resolution.
Its cohomological image in $H^{2r}_B$ is the filtration subspace
$F^0H^{2r}_B=E_\infty^{0,2r}$; there are no positive columns.
The quotient has only weights larger than $2r$. Frobenius
primary decomposition splits off exactly this weight-$2r$
subspace: its characteristic polynomial is coprime to that of
the quotient, since their Weil weights differ. This uses generalized
primary subspaces, not semisimplicity. Any equivariant map to
the pure weight-$2r$ target $H^{2r}(D_{\bar k})$ kills the
remaining subspace. Its image therefore comes entirely from
column zero, which purity identifies with the displayed individual
divisor Gysin maps. The reverse inclusion follows by factoring
each individual support map through the union. This proves
\eqref{r16gg:support-image}.

The primary ingredients are the same precise ones as in
Section~\ref{r16mb:section}: smooth-pair purity in SGA 4,
Expose XVI, Theorem 3.7 and Remark 3.10(a); Deligne,
\emph{La conjecture de Weil I}, Theorem 1.6; and localization
for closed supports, also recorded in Stacks Project Tag
\texttt{09XP}. The displayed spectral-sequence indices and
the general-degree weight deduction have been given here
rather than assumed from the degree-four statement.
\end{proof}

\begin{definition}[Scalar total boundary images]
\label{r16gg:definition}
Let $R$ be the valuation ring of a finite extension of $\mathbf Q_p$,
with finite residue field $k$. Let $\mathcal V/R$ be integral,
regular, projective and strictly semistable of relative dimension
$n\geq3$, with smooth geometrically integral generic fiber and
reduced special fiber
\[
                         \mathcal V_k=\bigcup_{i=1}^{s}D_i,
                              \qquad s\geq2.
\]
Assume the $D_i$ are smooth projective geometrically integral,
and all closed strata are smooth. Let the dual graph, with
an edge whenever $E_{ij}=D_i\cap D_j$ is nonempty, be connected.
The intersections may be disconnected and may have larger
native fields of constants. Fix one relatively ample $H$,
one $\ell\ne p$, and $1\leq r\leq n$.

For $j\ne i$, let $j_{ij,i}:E_{ij}\hookrightarrow D_i$.
The scalar total-image hypothesis in degree $r$ is
\[
 \begin{gathered}
 \dim_{\mathbf Q_\ell}M_i^r=1\quad\hbox{for every }i,
 \\
 M_i^r=\sum_{j\ne i}\operatorname{im}\left(
 H^{2r-2}((E_{ij})_{\bar k},\mathbf Q_\ell(r-1))
 \xrightarrow{j_{ij,i*}}
 H^{2r}((D_i)_{\bar k},\mathbf Q_\ell(r))\right).
 \end{gathered}
 \tag{GG16.5}\label{r16gg:total-image}
\]
This is a rank condition on the \emph{sum} of all boundary images
at each vertex. Rank one for each separate edge does not imply it.
\end{definition}

\begin{theorem}[Direct rational graph correction]
\label{r16gg:correction}
Under \eqref{r16gg:total-image}, every geometrically homologically
trivial rational codimension-$r$ generic class has a Chow extension
whose component restrictions are all geometrically homologically
trivial, and hence numerically trivial.

More explicitly, for any actual Chow extension $\Theta$, put
$\theta_i=i_i^!\Theta$, $H_i=H|_{D_i}$, and
\[
 \lambda_i=\deg_k(\theta_iH_i^{n-r}),\qquad
 w_{ij}=\deg_k\bigl((H|_{E_{ij}})^{n-1}\bigr)=w_{ji},
 \tag{GG16.6}\label{r16gg:weights}
\]
with $w_{ij}=0$ if $E_{ij}$ is empty and $w_{ii}=0$.
There is a rational solution, unique modulo a constant vector, of
\[
                  \sum_jw_{ij}(v_j-v_i)=\lambda_i
                  \qquad(1\leq i\leq s).
 \tag{GG16.7}\label{r16gg:laplacian}
\]
The extension
\[
                 \Theta^{\rm ad}
                    =\Theta-\sum_i v_i\,i_{i*}H_i^{r-1}
 \tag{GG16.8}\label{r16gg:extension}
\]
has the asserted zero restrictions. For this fixed initial
$\Theta$, its Chow class is independent of the constant
ambiguity in $v$.
\end{theorem}

\begin{proof}
Take a pure operational $K_0$ representative of the chosen
$\Theta$, by frozen \texttt{r12n:K0}. Its degree-$2r$ geometric
Chern class on the model restricts to
$a_i=\operatorname{cl}_\ell(\theta_i)$ and has zero geometric
generic restriction. Apply proper nearby-cycle comparison on
the entire model and then restrict its zero generic image to
\[
                  U_i=D_i\setminus\bigcup_{j\ne i}E_{ij}.
\]
Since the model is reduced strictly semistable, the structural
morphism is smooth along $U_i$. The nearby-cycle unit is an
isomorphism there, so $a_i|_{(U_i)_{\bar k}}=0$.
This is the proper-global-then-open argument of
Lemma~\ref{r16tb:numerical}; no nonproper specialization
isomorphism is asserted. Localization and
Lemma~\ref{r16gg:support}, with the uniform twist $(r)$ restored,
show that $a_i\in M_i^r$.

For every nonempty native edge, the rational algebraic class
\[
        j_{ij,i*}c_{1,\ell}(H|_{E_{ij}})^{r-1}\in M_i^r
\]
has pairing $w_{ij}>0$ with $c_{1,\ell}(H_i)^{n-r}$.
Thus this pairing is a nonzero functional on the one-dimensional
$M_i^r$. There is a unique vector $e_i\in M_i^r$ with that
pairing equal to one. It has a native rational algebraic
representative obtained by dividing any such edge class by
$w_{ij}$. In particular
\[
                   a_i=\lambda_i e_i,\qquad
 j_{ij,i*}c_{1,\ell}(H|_{E_{ij}})^{r-1}=w_{ij}e_i.
 \tag{GG16.9}\label{r16gg:normalized-lines}
\]
All $\lambda_i$ are rational degrees of the actual chosen cycles.
No rational lifting of an arbitrary $\ell$-adic class is assumed.

Frozen \texttt{r12n:weighted-trace} applied to the global top
class given by the pure representative times $c_1(H)^{n-r}$
gives
\[
                             \sum_i\lambda_i=0.
 \tag{GG16.10}\label{r16gg:trace-zero}
\]
Its generic factor is geometrically zero and the special
multiplicities are exactly one. Native degrees in
\eqref{r16gg:weights} include every constant-field degree
of every connected component of an intersection.

Let $L$ be the rational symmetric matrix
$(Lv)_i=\sum_jw_{ij}(v_j-v_i)$. Its column sums are zero,
and over $\mathbf R$ its energy is
\[
                 v^tLv=-\sum_{i<j}w_{ij}(v_i-v_j)^2.
 \tag{GG16.11}\label{r16gg:energy}
\]
All nonzero weights are positive. Connectedness implies
$\ker L=\mathbf Q(1,\ldots,1)$ over $\mathbf Q$ as well.
Its rank is $s-1$, so its image is exactly the rational
hyperplane of coordinate sum zero. Equation
\eqref{r16gg:trace-zero} therefore gives the rational solution
\eqref{r16gg:laplacian}, unique modulo constants. This is
ordinary rational matrix inversion; no real solution is
substituted for a rational Chow correction.

The principal total-fiber relation gives
\[
 c_1(N_{D_i/\mathcal V})=-\sum_{j\ne i}[E_{ij}]
                         \quad\hbox{in }\operatorname{CH}^1(D_i).
\]
Refined transverse restriction and self-intersection thus make
the component restriction of $\sum_jv_j i_{j*}H_j^{r-1}$ equal to
\[
              \sum_{j\ne i}(v_j-v_i)
                      j_{ij,i*}(H|_{E_{ij}})^{r-1}.
 \tag{GG16.12}\label{r16gg:vertical-restrictions}
\]
Triple and higher strata add no term to this Cartier-divisor
identity. On cohomology, \eqref{r16gg:normalized-lines} turns
it into $(Lv)_ie_i=\lambda_i e_i=a_i$.
Subtracting proves that every component Chow restriction of
\eqref{r16gg:extension} has zero geometric cycle class and
therefore is numerically zero. The generic Chow class is
unchanged. Changing $v$ by a constant changes the correction
by a scalar multiple of
$H^{r-1}\sum_i[D_i]=H^{r-1}\operatorname{div}(\pi)$,
which is zero in the model Chow group. This proves the final
assertion as well.
\end{proof}

\begin{corollary}[Native actual activation]
\label{r16gg:actual}
Let $X/K$ be smooth projective geometrically integral. Suppose
that after a finite extension $F/K$ there is a smooth projective
geometrically integral $V/F$, with a projective dominant morphism
$V\to X_F$, and a model satisfying
Definition~\ref{r16gg:definition} in degree $r$.
Then every geometrically homologically trivial rational
codimension-$r$ class on $X$ has zero native actual obstruction.
Every specified cycle in such a class has a strong Green form
on its exact support complement, with the equality on all
original compact actual tests. If the model satisfies the
total-image condition in all degrees, the target has full
strong G in every positive codimension.
\end{corollary}

\begin{proof}
Theorem~\ref{r16gg:correction} supplies the numerically
admissible model lift of the refined generic pullback class.
Frozen \texttt{r12n:K0} supplies its pure operational
virtual-bundle representative, including zero supported
degree on every integral vertical test of the required
dimension. Apply Corollary~\ref{r16G:numerical-witness}
and its exact-support bridge in the original $K$-theory.
The model viewed over $K^\circ$ retains the ramification
factor in special multiplicities and the residue factor
in native vertical degrees; both preserve zero. All finite
degrees in dominant actual Bott--Chern detection are complete
native degrees. The argument does not assume that the
dominant map is finite away from the chosen cycle, and
does not push its Green form analytically.
\end{proof}

\begin{remark}[A dual restriction test for the total-image hypothesis]
\label{r16gg:dual-test}
The dimension of $M_i^r$ is the rank of the single tuple-restriction map
\[
 H^{2n-2r}((D_i)_{\bar k},\mathbf Q_\ell(n-r))
 \longrightarrow
 \bigoplus_{j\ne i}H^{2n-2r}((E_{ij})_{\bar k},\mathbf Q_\ell(n-r)).
 \tag{GG16.13}\label{r16gg:dual-rank}
\]
Indeed Poincare duality gives perfect pairings on the smooth
projective $D_i$ and on every geometric component of the
intersections. Gysin--restriction adjunction identifies
\eqref{r16gg:dual-rank} with the transpose of the sum of Gysin
maps whose image is $M_i^r$. Their ranks are equal. Disconnected
intersections use the direct sum of their component pairings.
This gives a concrete full-cohomology test, as in the central
calculation of Proposition~\ref{r16sc:hypersurface-images}.
It is the rank of the tuple that must be one; separate rank-one
restriction maps can jointly have larger rank.
\end{remark}

\paragraph{Limits.}
The rank condition concerns the sum of all geometric boundary
Gysin images at each component. Individually one-dimensional
edge images can span a larger space and do not suffice.
Nor is this condition asserted stable under arbitrary model
blowups or alterations. The proof is for reduced strictly
semistable models only: at a multiple component the open
single-component locus need not be smooth over the trait,
so the nearby-cycle step cannot be imported unchanged.
No universal existence of these model witnesses or unconditional
general-G conclusion is claimed.
